\documentclass[11pt]{article}

\usepackage[margin=1in]{geometry}
\usepackage{amsmath,amssymb,amsfonts}
\usepackage{algorithm}
\usepackage{algorithmic}
\usepackage{booktabs}

\title{Solution Method for the 3-ESP, 3-NSP Multi-Provider Extension}
\author{}
\date{}

\begin{document}
\maketitle

\section{Overview}

This document describes the solution method implemented in the multi-provider
experiment scripts under \texttt{scripts/multi\_provider}.  The method extends
the original ESP--NSP Stackelberg framework to a market with three edge service
providers (ESPs), three network service providers (NSPs), and multiple users.
It is intended as an experimental extension rather than a replacement of the
main theoretical analysis in the paper.  The implementation follows the same
paper-facing idea: service providers set prices in Stage I, and users respond
in Stage II through offloading, provider-pair selection, and resource purchase
decisions.

Let
\[
    \mathcal E=\{1,2,3\}, \qquad \mathcal N=\{1,2,3\},
\]
denote the ESP and NSP sets, respectively.  Each ESP \(e\in\mathcal E\) owns
computation capacity \(F_e\), unit cost \(c^E_e\), and price \(p^E_e\).  Each
NSP \(n\in\mathcal N\) owns bandwidth capacity \(B_n\), unit cost \(c^N_n\),
and price \(p^N_n\).  Each user can either compute locally or select one
ESP--NSP pair \((e,n)\) for offloading.

\section{User Cost Model}

The implementation stores the user-side constants
\[
    a_i \triangleq \alpha_i w_i,
    \qquad
    \theta_{i,n}\triangleq
    \frac{d_i(\alpha_i+\beta_i\rho_i\varpi_i)}{\sigma_{i,n}},
\]
where \(\sigma_{i,n}\) is the spectral efficiency from user \(i\) to NSP
\(n\).  The local computation cost is denoted by \(C_i^\ell\).

If user \(i\) offloads through pair \((e,n)\), buys computation resource
\(f_i\), and buys bandwidth \(b_i\), its offloading cost is
\begin{equation}
    C^o_{i,e,n}(f_i,b_i)
    =
    \frac{a_i}{f_i}
    +
    \frac{\theta_{i,n}}{b_i}
    +
    p^E_e f_i
    +
    p^N_n b_i .
    \label{eq:multi-provider-offload-cost}
\end{equation}

The binary assignment variables can be written as
\[
    x_{i,e,n}\in\{0,1\},\qquad
    \sum_{e\in\mathcal E}\sum_{n\in\mathcal N}x_{i,e,n}\le 1.
\]
When \(\sum_{e,n}x_{i,e,n}=0\), user \(i\) computes locally.

\section{Stage-II User Response}

Given provider prices
\[
    \boldsymbol p=(p^E_1,p^E_2,p^E_3,p^N_1,p^N_2,p^N_3),
\]
the ideal Stage-II problem is a social-cost minimization over user assignment
and resource allocation:
\begin{align}
    \min_{\boldsymbol x,\boldsymbol f,\boldsymbol b}\quad
    &
    \sum_i
    \left(1-\sum_{e,n}x_{i,e,n}\right)C_i^\ell
    +
    \sum_i\sum_{e,n}x_{i,e,n}C^o_{i,e,n}(f_i,b_i)
    \label{eq:multi-provider-scm}
    \\
    \text{s.t.}\quad
    &
    \sum_i\sum_n x_{i,e,n}f_i \le F_e,
    \qquad \forall e\in\mathcal E,
    \nonumber\\
    &
    \sum_i\sum_e x_{i,e,n}b_i \le B_n,
    \qquad \forall n\in\mathcal N,
    \nonumber\\
    &
    C^o_{i,e,n}(f_i,b_i)\le C_i^\ell,
    \qquad \forall i,e,n \text{ with }x_{i,e,n}=1.
    \nonumber
\end{align}
This mixed-integer problem is expensive because users must choose among nine
provider pairs.  The experiment therefore uses a greedy assignment heuristic
with exact fixed-assignment resource response.

\subsection{Fixed-Assignment Closed-Form Response}

Consider a fixed assignment \(A\), where \(A_i=(e_i,n_i)\) for offloading users
and \(A_i=\emptyset\) for local users.  Let
\[
    \mathcal I_e(A)=\{i:A_i=(e,n)\text{ for some }n\},
    \qquad
    \mathcal I_n(A)=\{i:A_i=(e,n)\text{ for some }e\}.
\]
For each ESP and NSP, define the congestion threshold prices
\begin{equation}
    \bar p^E_e(A)
    =
    \left(
    \frac{\sum_{i\in\mathcal I_e(A)}\sqrt{a_i}}{F_e}
    \right)^2,
    \qquad
    \bar p^N_n(A)
    =
    \left(
    \frac{\sum_{i\in\mathcal I_n(A)}\sqrt{\theta_{i,n}}}{B_n}
    \right)^2 .
    \label{eq:multi-provider-bar-prices}
\end{equation}
The effective prices are
\begin{equation}
    \tilde p^E_e(A)=\max\{p^E_e,\bar p^E_e(A)\},
    \qquad
    \tilde p^N_n(A)=\max\{p^N_n,\bar p^N_n(A)\}.
    \label{eq:multi-provider-effective-prices}
\end{equation}
For any offloading user assigned to pair \((e,n)\), the fixed-assignment
resource response is
\begin{equation}
    f_i^*(A)=\sqrt{\frac{a_i}{\tilde p^E_e(A)}},
    \qquad
    b_i^*(A)=\sqrt{\frac{\theta_{i,n}}{\tilde p^N_n(A)}} .
    \label{eq:multi-provider-fixed-response}
\end{equation}
The corresponding capacity shadow prices are
\begin{equation}
    \lambda_e(A)=\tilde p^E_e(A)-p^E_e,
    \qquad
    \nu_n(A)=\tilde p^N_n(A)-p^N_n.
    \label{eq:multi-provider-shadow-prices}
\end{equation}
The implementation uses these expressions in
\texttt{fixed\_assignment\_response}.  They also make it cheap to evaluate the
social cost and individual-rationality margin of an assignment:
\begin{equation}
    m_i(A)=C_i^\ell-C^o_{i,e_i,n_i}\bigl(f_i^*(A),b_i^*(A)\bigr).
    \label{eq:multi-provider-margin}
\end{equation}
An assigned user is feasible only when \(m_i(A)\ge 0\).

\subsection{Greedy Admission and Rollback}

Starting from the empty offloading assignment, the Stage-II heuristic admits
users one by one.  Given the current assignment \(A\), each currently local
user \(j\) evaluates every ESP--NSP pair using the congestion-adjusted score
\begin{equation}
    h_{j,e,n}(A)
    =
    \min_{0<f\le F_e}
    \left\{
        \frac{a_j}{f}
        +(p^E_e+\lambda_e(A))f
    \right\}
    +
    \min_{0<b\le B_n}
    \left\{
        \frac{\theta_{j,n}}{b}
        +(p^N_n+\nu_n(A))b
    \right\}
    -
    C_j^\ell .
    \label{eq:multi-provider-score}
\end{equation}
Each one-dimensional minimization has the bounded closed form
\begin{equation}
    \min_{0<z\le U}\left\{\frac{\xi}{z}+\tau z\right\}
    =
    \begin{cases}
    2\sqrt{\xi\tau}, & \sqrt{\xi/\tau}\le U,\\[2mm]
    \xi/U+\tau U, & \sqrt{\xi/\tau}>U.
    \end{cases}
    \label{eq:bounded-1d-min}
\end{equation}
The algorithm selects the user-pair triple with the most negative score:
\[
    (j^*,e^*,n^*)
    \in
    \arg\min_{j,e,n} h_{j,e,n}(A).
\]
If the minimum score is nonnegative, the algorithm stops.  Otherwise, user
\(j^*\) is tentatively assigned to \((e^*,n^*)\), the fixed-assignment response
is recomputed, and a rollback check is applied.  The tentative admission is
removed if it makes the realized social cost worse or if any offloading user
violates individual rationality.  Users that fail the rollback check are
removed from the active candidate pool for that Stage-II solve.

The Stage-II output contains the final assignment, resources, shadow prices,
social cost, number of accepted admissions, number of rollbacks, and the social
cost trace.

\section{Provider Revenues}

After the Stage-II response is computed, each provider's realized demand and
revenue are evaluated from the assigned users:
\begin{equation}
    U^E_e(\boldsymbol p)
    =
    (p^E_e-c^E_e)
    \sum_{i\in\mathcal I_e(A^*(\boldsymbol p))} f_i^*,
    \qquad e\in\mathcal E,
    \label{eq:esp-revenue}
\end{equation}
\begin{equation}
    U^N_n(\boldsymbol p)
    =
    (p^N_n-c^N_n)
    \sum_{i\in\mathcal I_n(A^*(\boldsymbol p))} b_i^*,
    \qquad n\in\mathcal N.
    \label{eq:nsp-revenue}
\end{equation}
Here \(A^*(\boldsymbol p)\) denotes the assignment returned by the Stage-II
heuristic at prices \(\boldsymbol p\).

\section{Stage-I Six-Provider Pricing Dynamics}

The Stage-I game has six leaders: three ESPs and three NSPs.  Exhaustive
six-dimensional price search is not practical.  The implementation therefore
uses a boundary-price best-response heuristic inspired by the switching-boundary
logic of the main paper.

\subsection{Local Candidate Assignment Family}

At a current price vector \(\boldsymbol p\), the solver first computes the
Stage-II response \(A^*(\boldsymbol p)\).  It then constructs a local family of
candidate assignments.  Current offloading users are sorted by increasing
individual-rationality margin \(m_i\), so the first users are the easiest to
remove.  Current local users are sorted by their best pair score
\(\min_{e,n}h_{i,e,n}\), so the first users are the most likely to enter.

For a small integer \(Q\), candidate assignments are created by removing the
first \(r\) fragile offloading users and adding the first \(s\) promising local
users, where
\[
    0\le r\le Q,\qquad 0\le s\le Q.
\]
Duplicate and empty assignments are discarded.  Hence the candidate family is
small, with size at most \((Q+1)^2\) before duplicate removal.

\subsection{Boundary Price for One Provider}

For a provider \(k\), all other providers' prices are fixed.  For each local
candidate assignment \(A\), the solver estimates the right boundary of the
assignment's feasible price interval for provider \(k\).  This boundary is the
first price at which one user served by provider \(k\) becomes indifferent
between offloading and local computation:
\begin{equation}
    \min_{i\in\mathcal I_k(A)}
    m_i(A;p_k,p_{-k})=0.
    \label{eq:multi-provider-boundary-price}
\end{equation}
In the implementation this boundary is computed by bracketing and bisection.
If the candidate assignment is already infeasible at the provider's cost floor,
it is skipped.  If no zero-margin crossing is found before the configured
maximum price, the maximum reachable price is used as the candidate boundary.

For each provider \(k\), the evaluated price set contains:
\begin{itemize}
    \item the current price \(p_k\);
    \item all valid boundary prices from the local assignment family;
    \item one small upward exploratory price, \(\min\{p_k^{\max},\max(1.1p_k,p_k+0.03)\}\).
\end{itemize}
Each price in this finite set is evaluated by rerunning the Stage-II heuristic
under the corresponding unilateral deviation.  The provider's estimated best
response is the evaluated price with the largest realized revenue:
\begin{equation}
    \widehat p_k^{\mathrm{BR}}
    \in
    \arg\max_{q_k\in\widehat{\mathcal C}_k(\boldsymbol p)}
    U_k(q_k,p_{-k}).
    \label{eq:multi-provider-br}
\end{equation}
The restricted best-response gain is
\begin{equation}
    \widehat G_k(\boldsymbol p)
    =
    U_k(\widehat p_k^{\mathrm{BR}},p_{-k})-U_k(p_k,p_{-k}).
    \label{eq:multi-provider-gain}
\end{equation}

\subsection{Iterative Pricing Rule}

The restricted Nash-equilibrium gap is defined as
\begin{equation}
    \widehat\delta(\boldsymbol p)
    =
    \max_{k\in\mathcal E\cup\mathcal N}\widehat G_k(\boldsymbol p).
    \label{eq:multi-provider-restricted-gap}
\end{equation}
At each Stage-I iteration, all six providers compute their estimated
best-response gains.  If
\[
    \widehat\delta(\boldsymbol p)\le \varepsilon,
\]
the pricing dynamics terminates.  Otherwise, the implementation updates the
single provider with the largest restricted gain:
\[
    k^* \in \arg\max_k \widehat G_k(\boldsymbol p),
    \qquad
    p_{k^*}\leftarrow \widehat p_{k^*}^{\mathrm{BR}}.
\]
All other provider prices remain unchanged in that iteration.  After the final
Stage-I iteration, the solver runs Stage II one more time and reports the final
assignment, provider revenues, restricted gap, and resource utilization.

\section{Algorithm Summary}

\begin{algorithm}[h]
\caption{Multi-Provider Stage-II Heuristic}
\begin{algorithmic}[1]
\REQUIRE Prices \(\boldsymbol p\), capacities \(\{F_e\}\), \(\{B_n\}\)
\STATE Initialize empty assignment \(A\) and active local-user pool.
\REPEAT
    \STATE Compute fixed-assignment response by
    \eqref{eq:multi-provider-effective-prices}--\eqref{eq:multi-provider-fixed-response}.
    \STATE Roll back the last admitted user if the realized social cost does not decrease.
    \STATE Remove any assigned user violating \(m_i(A)\ge 0\).
    \STATE For every active local user and every pair \((e,n)\), compute \(h_{i,e,n}(A)\).
    \IF{the minimum score is negative}
        \STATE Tentatively assign the corresponding user to the corresponding pair.
    \ELSE
        \STATE Stop.
    \ENDIF
\UNTIL{the active pool is empty or the iteration limit is reached}
\STATE Recompute the final fixed-assignment response and remove infeasible users.
\RETURN assignment, resources, shadow prices, social cost, and diagnostics.
\end{algorithmic}
\end{algorithm}

\begin{algorithm}[h]
\caption{Six-Provider Boundary-Price Pricing Dynamics}
\begin{algorithmic}[1]
\REQUIRE Initial prices \(\boldsymbol p^{(0)}\), local family size \(Q\), tolerance \(\varepsilon\)
\FOR{\(t=0,1,\ldots,T_{\max}-1\)}
    \STATE Run Stage II at \(\boldsymbol p^{(t)}\).
    \FORALL{providers \(k\in\mathcal E\cup\mathcal N\)}
        \STATE Construct the local candidate assignment family.
        \STATE Compute candidate boundary prices using \eqref{eq:multi-provider-boundary-price}.
        \STATE Evaluate each candidate unilateral price by rerunning Stage II.
        \STATE Obtain \(\widehat p_k^{\mathrm{BR}}\) and \(\widehat G_k\).
    \ENDFOR
    \STATE Set \(\widehat\delta(\boldsymbol p^{(t)})=\max_k\widehat G_k\).
    \IF{\(\widehat\delta(\boldsymbol p^{(t)})\le\varepsilon\)}
        \STATE Stop.
    \ENDIF
    \STATE Update only the provider \(k^*\in\arg\max_k\widehat G_k\).
\ENDFOR
\STATE Run Stage II at the final price vector.
\RETURN final prices, assignment, revenues, utilization, and diagnostics.
\end{algorithmic}
\end{algorithm}

\section{Implementation Parameters Used in the Reported Runs}

The reported \(N=30\) and \(N=50\) runs use the
\texttt{code\_default\_heterogeneous} setup, which follows the numerical scale
of the experiment code rather than the physical-unit scale in the paper text.
The user distributions are
\[
\begin{array}{lll}
w_i\sim U[0.5,2.5], & d_i\sim U[1.0,10.0], &
f_i^\ell\sim U[0.5,1.2],\\
\alpha_i\sim U[10.0,15.0], &
\beta_i\sim U[0.1,0.5], &
\rho_i\sim U[0.5,2.0],\\
\varpi_i\sim U[0.8,1.2], &
\kappa_i\sim U[0.01,0.05], &
\sigma^{\mathrm{base}}_{i,n}\sim U[0.5,3.0].
\end{array}
\]
The NSP channel qualities use the bias vector
\[
    [1.15,\;1.00,\;0.85],
\]
clipped to the interval \([0.5,3.0]\).  The provider parameters are
\[
    (F_1,F_2,F_3)=(45,35,25),
    \qquad
    (B_1,B_2,B_3)=(20,14,10),
\]
\[
    (c^E_1,c^E_2,c^E_3)=(0.008,0.010,0.013),
    \qquad
    (c^N_1,c^N_2,c^N_3)=(0.007,0.010,0.014).
\]
The default initial prices are
\[
    \boldsymbol p_E^{(0)}=(0.42,0.50,0.58),
    \qquad
    \boldsymbol p_N^{(0)}=(0.40,0.50,0.62).
\]
The \(N=30\) and \(N=50\) runs used \(Q=2\), maximum price \(6.0\) for both
resource types, ten Stage-I iterations, and ninety-six Stage-II iterations per
Stage-II solve.

\section{Remarks}

The method is intentionally heuristic.  It preserves the main paper's
boundary-price intuition, but the full 3-ESP, 3-NSP setting has a
six-dimensional leader game and a nine-choice assignment problem on the user
side.  Therefore, the computed outcome should be interpreted as an approximate
multi-provider Stackelberg outcome produced by a boundary-price restricted
best-response dynamics.  The diagnostics saved by the scripts, including
assignment heatmaps, restricted-gap trajectories, provider revenues, and
resource utilizations, are used to judge whether the resulting outcome is
economically meaningful.

\end{document}
